% !TEX root = ../matan.tex

\section{Длина гладкого пути}

Напомним, что кривая $\Gamma$ называется \textbf{спрямляемой}, если множество
длин вписанных в неё ломаных ограничено сверху, и тогда
\[
l(\Gamma):=\sup\{\,l(\Gamma')\st \Gamma'\ \text{--- ломаная, вписанная в }\Gamma\,\}.
\]
Длина звена ломаной с вершинами $\vphi(t_j),\vphi(t_{j+1})$ равна
\[\bigl|\vphi(t_{j+1})-\vphi(t_j)\bigr|
=\sqrt{\sum_{n=1}^d\bigl(\vphi_n(t_{j+1})-\vphi_n(t_j)\bigr)^2}.\]

\begin{Theorem}[required]{Длина гладкого пути}{}
	Пусть $\Gamma$ --- \textbf{гладкая} кривая в $\RR^d$, то есть
	\[
	\Gamma=\vphi(I),\qquad I=[a,b],\qquad
	\vphi=(\vphi_1,\ldots,\vphi_d),\qquad \vphi_n:I\to\RR,
	\]
	и все $\vphi_n'\in C([a,b])$ (на концах --- односторонние производные).
	Тогда $\Gamma$ спрямляема, и
	\[
	l(\Gamma)=\int_a^b \sqrt{\sum_{n=1}^d \bigl(\vphi_n'(t)\bigr)^2}\,dt.
	\]
\end{Theorem}

\begin{proof}
	Положим
	\[
	L:=\int_a^b \sqrt{\sum_{n=1}^d (\vphi_n'(t))^2}\,dt.
	\]
	Докажем два неравенства: $l(\Gamma')\le L$ для любой вписанной ломаной
	(откуда $l(\Gamma)=\sup l(\Gamma')\le L$) и существование ломаных
	$\Gamma_k'$ с $l(\Gamma_k')\to L$ (откуда $l(\Gamma)\ge L$). Из них следует
	$l(\Gamma)=L$.
	
	\medskip
	\textbf{1. Оценка сверху:} $\forall\,\Gamma'\ \ l(\Gamma')\le L$.
	
	Пусть $\{t_j\}_{j=0}^N$ --- точки параметризации ломаной. Тогда
	\[
	l(\Gamma')=\sum_{j=0}^{N-1}
	\sqrt{\sum_{n=1}^d \bigl(\vphi_n(t_{j+1})-\vphi_n(t_j)\bigr)^2}.
	\]
	Достаточно показать, что для каждого $j$
	\[
	\sqrt{\sum_{n=1}^d \bigl(\vphi_n(t_{j+1})-\vphi_n(t_j)\bigr)^2}
	\le
	\int_{t_j}^{t_{j+1}} \sqrt{\sum_{n=1}^d (\vphi_n'(t))^2}\,dt;
	\]
	суммируя по $j$, получим $l(\Gamma')\le\int_a^b\sqrt{\sum_n(\vphi_n')^2}\,dt=L$.
	
	По формуле Ньютона--Лейбница
	$\vphi_n(t_{j+1})-\vphi_n(t_j)=\int_{t_j}^{t_{j+1}}\vphi_n'(t)\,dt$, поэтому
	достаточно проверить
	\[
	\sum_{n=1}^d \left(\int_{t_j}^{t_{j+1}} \vphi_n'(t)\,dt\right)^2
	\le
	\left(\int_{t_j}^{t_{j+1}} \sqrt{\sum_{n=1}^d (\vphi_n'(t))^2}\,dt\right)^2.
	\]
	Положим $f_n(t):=\bigl(\vphi_n'(t)\bigr)^2\ge0$. Тогда правая часть есть
	$\bigl(\int_{t_j}^{t_{j+1}}(\sum_n f_n)^{1/2}\,dt\bigr)^2$, а левая
	оценивается так:
	\[
	\sum_{n=1}^d \left(\int_{t_j}^{t_{j+1}} \vphi_n'(t)\,dt\right)^2
	\le
	\sum_{n=1}^d \left(\int_{t_j}^{t_{j+1}} |\vphi_n'(t)|\,dt\right)^2
	=
	\sum_{n=1}^d \left(\int_{t_j}^{t_{j+1}} (f_n(t))^{1/2}\,dt\right)^2.
	\]
	Теперь применяем \textbf{обратное неравенство Минковского при $p=\frac12$}
	и получаем
	\[
	\sum_{n=1}^d \left(\int_{t_j}^{t_{j+1}} (f_n(t))^{1/2}\,dt\right)^2
	\le
	\left(\int_{t_j}^{t_{j+1}} \Bigl(\sum_{n=1}^d f_n(t)\Bigr)^{1/2}\,dt\right)^2.
	\]

	Тем самым требуемая оценка доказана, и $l(\Gamma)=\sup l(\Gamma')\le L$.
	
	\medskip
	\textbf{2. Построение ломаных с $l(\Gamma_k')\to L$.}
	
	Возьмём диадическое дробление $\{t_{k,j}\}_{j=0}^{N_k}$ отрезка $[a,b]$ ранга
	$k$ и впишем в $\Gamma$ ломаную $\Gamma_k'=\{\vphi(t_{k,j})\}_{j=0}^{N_k}$,
	соединяя соседние точки $\vphi(t_{k,j}),\vphi(t_{k,j+1})$. По построению
	$l(\Gamma_k')\le l(\Gamma)$ для всех $k$. Покажем, что $l(\Gamma_k')\to L$.
	
	Функции $\vphi_n'$ непрерывны на компакте $[a,b]$, значит, равномерно
	непрерывны: для любого $\eps>0$ найдётся $k_0$ такое, что при $k\ge k_0$ для
	любого отрезка дробления $[t_{k,j},t_{k,j+1}]$ и любых $\tau_1,\tau_2$ из него
	\[
	|\vphi_n'(\tau_1)-\vphi_n'(\tau_2)|<\eps,\qquad n=1,\ldots,d.
	\]
	Записав $|t_{k,j+1}-t_{k,j}|=\int_{t_{k,j}}^{t_{k,j+1}}1\,dt$, представим
	разность в виде
	\[
	L-l(\Gamma_k')
	=\sum_{j=0}^{N_k-1}\int_{t_{k,j}}^{t_{k,j+1}}
	\left(\sqrt{\sum_{n=1}^d (\vphi_n'(t))^2}
	-\sqrt{\sum_{n=1}^d
		\frac{\bigl(\vphi_n(t_{k,j+1})-\vphi_n(t_{k,j})\bigr)^2}{|t_{k,j+1}-t_{k,j}|^2}}
	\right)dt.
	\]
	Для евклидовой нормы $\bigl|\sqrt{\sum_n\alpha_n^2}-\sqrt{\sum_n\beta_n^2}\bigr|
	\le\sqrt{\sum_n(\alpha_n-\beta_n)^2}$, поэтому
	\[
	L-l(\Gamma_k')\le\sum_{j=0}^{N_k-1}\int_{t_{k,j}}^{t_{k,j+1}}
	\sqrt{\sum_{n=1}^d\left(|\vphi_n'(t)|
		-\frac{\vphi_n(t_{k,j+1})-\vphi_n(t_{k,j})}{|t_{k,j+1}-t_{k,j}|}\right)^2}\,dt.
	\]
	По теореме Лагранжа о конечных приращениях для каждого $n$ найдётся точка
	$\tau_{k,j}^n\in[t_{k,j},t_{k,j+1}]$ такая, что
	$\dfrac{\vphi_n(t_{k,j+1})-\vphi_n(t_{k,j})}{|t_{k,j+1}-t_{k,j}|}
	=\vphi_n'(\tau_{k,j}^n)$. Поэтому
	\[
	L-l(\Gamma_k')\le\sum_{j=0}^{N_k-1}\int_{t_{k,j}}^{t_{k,j+1}}
	\sqrt{\sum_{n=1}^d\bigl(\vphi_n'(t)-\vphi_n'(\tau_{k,j}^n)\bigr)^2}\,dt.
	\]
	При $k\ge k_0$ имеем $|\vphi_n'(t)-\vphi_n'(\tau_{k,j}^n)|<\eps$, откуда
	\[
	L-l(\Gamma_k')\le\sum_{j=0}^{N_k-1}\int_{t_{k,j}}^{t_{k,j+1}}
	\sqrt{\sum_{n=1}^d\eps^2}\,dt
	=\eps\sqrt{d}\sum_{j=0}^{N_k-1}|t_{k,j+1}-t_{k,j}|
	=\eps\sqrt{d}\,(b-a).
	\]
	Итак, $0\le L-l(\Gamma_k')\le\eps\sqrt{d}\,(b-a)$ при $k\ge k_0$. В силу
	произвольности $\eps$ получаем $l(\Gamma_k')\to L$, а значит, $l(\Gamma)\ge L$.
	
	\medskip
	Из пунктов 1 и 2: $L\le l(\Gamma)\le L$, то есть $\Gamma$ спрямляема и
	\[
	l(\Gamma)=\int_a^b \sqrt{\sum_{n=1}^d (\vphi_n'(t))^2}\,dt. \qedhere
	\]
\end{proof}

\begin{Example}{Параметризация окружности}{}
	Для окружности радиуса $r>0$:
	\[
	\vphi_1(t)=r\cos(2\pi t),\qquad \vphi_2(t)=r\sin(2\pi t),\qquad t\in[0,1].
	\]
	Тогда $\vphi_1'=-2\pi r\sin(2\pi t)$, $\vphi_2'=2\pi r\cos(2\pi t)$, и
	\[
	l(\Gamma)=\int_0^1\sqrt{r^2(2\pi)^2\bigl(\sin^2(2\pi t)+\cos^2(2\pi t)\bigr)}\,dt
	=\int_0^1 2\pi r\,dt=2\pi r.
	\]
\end{Example}